% ====================================================================
%  GAYA BERSAMA MODUL OLIMPIADE (dipakai teobil_isi.tex & kombin_isi.tex)
%  Semua tampilan hitam-putih / abu-abu supaya hemat tinta.
% ====================================================================
\usepackage[bahasai]{babel}
\usepackage{amsmath, amssymb}
\usepackage{geometry}
\geometry{top=2.5cm, bottom=2.5cm, left=2.5cm, right=2.5cm}
\usepackage[most]{tcolorbox}
\usepackage{enumitem}
\usepackage{array}
\usepackage{comment}
\usepackage{tikz}

\setlength{\parindent}{0pt}
\setlength{\parskip}{5pt}

\tcbset{colback=white, colframe=black, colbacktitle=white, coltitle=black,
  fonttitle=\bfseries, boxrule=0.6pt, arc=1mm,
  before upper={\setlength{\parskip}{5pt}}}

% Kegiatan pembuka
\newtcolorbox{selidiki}{enhanced, breakable, colback=black!4, colframe=black!60,
  colbacktitle=black!15, coltitle=black, title={Ayo Selidiki!}}

% Contoh soal (abu-abu tipis, garis tebal di kiri)
\newtcolorbox{contoh}{enhanced, breakable, colback=black!4, colframe=black!55,
  boxrule=0pt, leftrule=3pt, arc=0pt, left=3mm, before skip=6pt, after skip=6pt}

% Soal untuk dicoba langsung (putih, garis putus-putus)
\newtcolorbox{giliranmu}{enhanced, borderline={0.6pt}{0pt}{black, dashed},
  frame hidden, title={Sekarang Giliranmu!},
  attach boxed title to top left={xshift=3mm, yshift=-2mm},
  boxed title style={colframe=white, colback=white}}

% Ringkasan akhir topik
\newtcolorbox{rangkuman}{enhanced, colback=white, colframe=black!60,
  colbacktitle=black!15, coltitle=black, title={Rangkuman}}

% Garis isian
\newcommand{\kosong}{\rule{1.2cm}{0.4pt}}

% Label level: blok abu-abu kecil
\newcommand{\level}[1]{\par\medskip\noindent\colorbox{black!12}{\strut\textbf{#1}}\par\smallskip}

% Pembahasan soal: hanya muncul di Edisi Guru
\newtcolorbox{kotakjawab}{enhanced, breakable, boxrule=0pt, leftrule=2pt, arc=0pt,
  colback=white, colframe=black, left=2mm, top=1mm, bottom=1mm,
  before skip=3pt, after skip=3pt}
\newcommand{\jawab}[2]{\ifguru
  \begin{kotakjawab}\small\textbf{Jawaban: #1.} #2\end{kotakjawab}\fi}

% Catatan mengajar: hanya muncul di Edisi Guru
\ifguru
  \newtcolorbox{catatanguru}{breakable, colback=white, colframe=black!60,
    colbacktitle=black!15, title={Catatan untuk Guru}}
\else
  \excludecomment{catatanguru}
\fi

\newcommand{\C}[2]{\binom{#1}{#2}}
